\section{Introduction}\label{sec:introduction}

Let $F=\mathbb Q(\sqrt{-3})$, let $\mathcal O$ be its ring of integers,
and write $N\mathfrak a=|\mathcal O/\mathfrak a|$ for the norm of a
nonzero integral ideal. A finite-order Hecke character $\eta$ modulo
a nonzero integral ideal $\mathfrak f$ is a character of the ray class group
modulo $\mathfrak f$, extended by zero to ideals not coprime to $\mathfrak f$.
This agrees with the usual finite-order idelic definition: the complex
place contributes no nontrivial finite-order continuous character
\cite[Chapter VI, Section 1]{MilneCFT}. For $\Re s>1$, put
\[
 L_F(s,\eta)=\sum_{\mathfrak a\ne0}
       \frac{\eta(\mathfrak a)}{(N\mathfrak a)^s}
 =\prod_{\mathfrak p}
       \left(1-\eta(\mathfrak p)(N\mathfrak p)^{-s}\right)^{-1}.
\]
The same notation denotes its meromorphic continuation. For a
Dirichlet character $\chi$ modulo $q$, extended by zero on nonunits, we
likewise write
\[
 L(s,\chi)=\sum_{n\ge1}\frac{\chi(n)}{n^s}
 \qquad(\Re s>1)
\]
and then continue meromorphically; the character modulo $1$ gives the
Riemann zeta function $\zeta(s)$. The problem here is uniform exclusion:
to find a constant $\sigma_0<1$, independent of the character, its
conductor, and the height, such that these functions have no zeros in
$\Re s>\sigma_0$.

For $\zeta$ alone, the existence of a fixed $\sigma_0<1$ for which
$\zeta(s)\ne0$ in $\Re s>\sigma_0$ is called the quasi-Riemann
hypothesis \cite[Section 1]{BCSS}. It asks for one gap valid at every
height, not merely nonvanishing on $\Re s=1$. The corresponding uniform
formulation over all Dirichlet characters appears in
\cite[p.~288, Equation (1.4) and the following paragraph]{FG}.

The connection between these functions and prime distribution has a
long history. Dirichlet's 1837 proof of infinitude of primes in every
reduced arithmetic progression introduced the character $L$-series
that separate residue classes \cite{Dirichlet}. Riemann's 1859 memoir
related the zeros of $\zeta$ to the distribution of primes and formulated
the critical-line conjecture \cite{Riemann}. Hadamard and de la
Vall\'ee Poussin independently proved in 1896 that $\zeta$ has no zero
on $\Re s=1$, obtaining the prime number theorem
\cite{Hadamard,ValleePoussin}. Hecke subsequently developed the
character zeta functions and their analytic theory over number fields
\cite{Hecke}.

Classical zero-free regions for Hecke $L$-functions approach the line
$\Re s=1$ as the conductor or height increases and can allow one
simple real exceptional zero; a precise form for abelian extensions is given
by \cite[Theorem 3.1]{TZ}. Such regions do not give uniform exclusion in
a fixed half-plane. Zero-density estimates address a different question:
they bound the number of zeros to the right of a given vertical line.
For example, Guth and Maynard's large-value estimates improve zero-density
bounds for $\zeta$ \cite[Theorem 1.2]{GM}, but these zero-density bounds
do not exclude every zero.

\begin{theorem}\label{thm:main}
Every finite-order Hecke $L$-function over
$F=\mathbb Q(\sqrt{-3})$ has no zero in $\Re s>7/8$.
The same holds for every Dirichlet $L$-function, including $\zeta(s)$.
A pole at $s=1$ for a principal character is allowed.
\end{theorem}

In particular, Theorem~\ref{thm:main} resolves the quasi-Riemann
hypothesis affirmatively: the supremum of the real parts of the
nontrivial zeros of $\zeta$, namely its zeros in $0<\Re s<1$, is at most
$7/8$. The boundary $\Re s=7/8$ is not included, while any exceptional
real zero in $(7/8,1)$ is excluded. The theorem does not establish the
Riemann hypothesis or its generalized versions, which place nontrivial
zeros on $\Re s=1/2$. The Riemann hypothesis remains open \cite{Clay}.

Kubota's metaplectic theory and Patterson's cubic theta series provide
the automorphic setting \cite{Kubota,Patterson}. We use the unconditional
explicit cusp expansions recorded by Dunn and Radziwi\l\l{}
\cite[Section 5 and Appendix A]{DR}; their GRH-conditional prime asymptotic
is not an input. The corresponding first-moment asymptotic is proved
unconditionally in an independent manuscript
\cite[Theorem~1.1]{OpenAICubicFirstMoment}; that result is also not used
here. Proposition~\ref{prop:completed-reflection} derives the
reflection used here while retaining the characters' zero-on-nonunit
restrictions.

The character large-sieve arguments build on the quadratic Hecke-family
estimate of Goldmakher and Louvel \cite[Definition 1 and Theorem 1.1]{GL},
the higher-order norm recursion of Blomer, Goldmakher, and Louvel
\cite[Theorem 1.3 and Section 3]{BGL}, and Heath-Brown's cubic estimate
\cite[Theorem 2]{HB}. We prove the precise sextic specialization needed
here in Lemma~\ref{lem:sextic-large-sieve}, using paired residue characters
to carry out that recursion in the primary-generator convention. An
Eisenstein-integer form is also recorded by Gao and Zhao
\cite[Lemma 2.9]{GZSextic}. The recursive
moment arguments are related to the framework of Heath-Brown's quadratic
method \cite[Section 2]{HBQuadratic}. We also use the ordinary Hecke
functional equation \cite[Equation (1.1)]{GZ} and prime counting in a
fixed ray class \cite[Theorem 1.1]{TZ}.

The planar additive large sieve used in the first stage is classical;
compare Huxley's multivariable and number-field inequality \cite{Huxley}
and the Poisson proof in \cite[Section 3, Theorem 3]{BaierBansal}.
We include a direct Eisenstein-lattice proof to record the normalization
needed for the reduced-fraction expansion. David, de Faveri, Dunn, and
Stucky combine Patterson's coefficients, the cubic large sieve, and
mollified moments to prove nonvanishing at $s=1/2$ for a positive
proportion of a cubic Hecke family
\cite[Theorem 1.1 and Section 1.2]{DFDS}. The zero detector below uses
the classical truncated-inverse mechanism; compare
\cite[Appendix C]{MP}, and for a higher-order-character
density application \cite[Corollary 1.6 and Section 5]{BGL}.

The proof has two stages, each comparing two representations of a
completed cubic-theta sum but using a different normalized sum in the
continuation argument. Part~\ref{part:eleven-twelfths} proves the corresponding
$11/12$ assertion in Theorem~\ref{thm:eleven-twelfths}, already obtaining a
fixed zero-free half-plane for both families.
Part~\ref{part:seven-eighths} starts from that conclusion and introduces
prime compensation, asymmetric scales, and two additional moment estimates.
The contribution developed here is the construction of these compatible
reflected and Poisson comparisons, including the principal residues, with
positive exponent margins chosen independently of the target character.

\subsection*{Proof overview}

\paragraph{A common continuation principle.}
The analytic argument begins with the family of primitive finite-order
Hecke characters over $F$. Let $\beta_*$ be the supremum of $1/2$ and the
real parts of their zeros in $1/2\le\Re s\le1$; poles are not included.
Imprimitive characters have the same possible zeros in $\Re s>0$, because
the finitely omitted Euler factors are nonzero there. If $\beta_*$ exceeds
a proposed boundary $\sigma_0$, the task is to continue every target
reciprocal across a common positive distance to the left of $\beta_*$.

Section~\ref{sec:analytic-reduction} gives the precise criterion. For each
target $\eta$ and large real scale $Z$, it compares a normalized character
sum with a Mellin integral containing $1/L_F(s,\eta)$, after deletion of
finitely many Euler factors and multiplication by a holomorphic factor
bounded away from zero. A direct bound for the sum and a power-saving bound
for its difference from that integral imply the required continuation.
The two parts use this same principle with different affine powers of $Z$
in the Mellin integral:
\[
 C_{\mathrm I}(s)=s-\frac23,\qquad
 C_{\mathrm{II}}(s)=s-\frac{11}{16}.
\]
Thus the common analytic principle does not identify the two normalized
sums. Only the positive power margins must be independent of $\eta$;
fixed-character constants and lower thresholds may depend on it.

\paragraph{The completed sum.}
Section~\ref{sec:arithmetic} sets up the sextic residue characters over
$\mathcal O$, always retaining their zero values on nonunits. The base
sum in Section~\ref{sec:probe} averages smoothed cubic-theta Fourier
coefficients against these characters and a fixed target $\eta$.
After the fixed rescaling in the theta expansion, the completed indices
are $cn^3$, where $c,n\in\mathcal O$ are congruent to $1$ modulo $3$ and
$c$ is squarefree. The variables $c$ and $n$ may share prime factors.
Here completion means retaining the full cubic factor $n^3$, rather than
restricting the index to its squarefree part.
The target character and a finite ray-class phase act on the whole product
$cn^3$, rather than separately on its two factors.

The resulting completed base sum has two exact representations.
The reflection in Proposition~\ref{prop:completed-reflection}
transforms its theta coefficients while retaining the character zeros in
the resulting formula.
Poisson summation transforms the averaging variable. Its nonzero frequencies
have the form $ua^6$, where $u,a\in\mathcal O$ and $u$ is sixth-power-free:
every prime valuation of $u$ is at most five.
Section~\ref{sec:local-euler}
then expresses the contribution of each such row through a quotient of
Hecke $L$-functions and a controlled Euler product. For $u=1$, this
quotient contains the reciprocal of the target function.

\paragraph{The balanced first stage.}
In Part~\ref{part:eleven-twelfths}, the two averaging scales are both
$Z^{1/2}$. On the reflected side, the base sum separates into a completed
theta row and an additive polynomial. The quadratic large sieve bounds
the mean square of the reflected rows. Expanding the additive polynomial
produces reduced fractions in $\mathbb C/\mathcal O$; their separation and
coefficient mass give the required planar large-sieve bound. The
Cauchy--Schwarz inequality combines the two norms in
Proposition~\ref{prop:balanced-low}.
This direct estimate does not use the later inverse-moment recursion.

On the Poisson side, the intermediate rows are grouped by their norm and by
the location of zeros in bounded-height rectangles for the finite family
of Hecke twists that each row determines. Section~\ref{sec:detector} assigns
zero-free rectangles to these families. When a row has a selected zero
above the detector's fixed floor for the real part, it produces two large
Dirichlet polynomials: one with
ideal M\"obius coefficients, representing a truncated reciprocal, and one
without those coefficients. They have one common row character and one
common twist height. Part~\ref{part:eleven-twelfths} uses only the inverse
polynomial in its row count. After the part of the row with prime valuations
at least two is fixed, the sextic large sieve applies to its squarefree
factor; Proposition~\ref{prop:sextic-row-count} gives the resulting count.
Rows at the floor and the remaining small and large norm ranges are bounded
directly.

The principal row is treated separately. Its residues, together with the
local Euler identity, give a nonzero scalar multiple of the target Mellin
integral. After normalization, the direct estimate and the remaining-row
estimate verify the common continuation criterion with
$C_{\mathrm I}(s)=s-2/3$. This proves
Theorem~\ref{thm:eleven-twelfths}. The entire family is needed even for the
consequence about $\zeta$, because the Poisson rows introduce finite-order
Hecke twists of the target.

\paragraph{The refined second stage.}
Part~\ref{part:seven-eighths} retains the completed support and the shared
arithmetic identities, but it modifies the base sum. Selected prime factors
provide a local compensation that cancels an unwanted Euler contribution,
and the two averaging scales are no longer equal. After the new residue
calculation and scalar normalization, the corresponding Mellin signal uses
$C_{\mathrm{II}}(s)=s-11/16$. The low estimate for the normalized sum again
follows directly from reflected energy and an additive mean-square estimate;
Section~\ref{sec:compensated-low} proves it for the modified sum.

The high estimate uses both polynomials supplied by the zero detector.
For a smooth compactly supported $W$ on $(0,\infty)$, a scale $D>0$,
and a finite-order Hecke character $\psi$, their basic forms are
\[
 D^{-1/2}\sum_{\mathfrak a\ne0}\mu(\mathfrak a)\psi(\mathfrak a)
       W(N\mathfrak a/D),\qquad
 D^{-1/2}\sum_{\mathfrak a\ne0}\psi(\mathfrak a)W(N\mathfrak a/D),
\]
where $\mu$ is the ideal M\"obius function. In the specified intermediate
norm ranges above the detector floor, each retained row has a large inverse polynomial and a large plain
polynomial with one common row character and twist height. Bounds for
their moments limit the number of such rows. Section~\ref{sec:refined-row-counts}
combines the two bounds, using integer powers and selected prime factors
in the larger intermediate norm ranges. Small and large norm ranges,
and the floor class, remain direct estimates.

The two moment bounds require separate inductions.
Section~\ref{sec:inverse} proves the second-moment estimate for an inverse
polynomial multiplied by sums over disjoint prime sets
(Lemma~\ref{lem:marked}). Its recursive step uses two finite Poisson
transformations to shorten the row range, with reflected energy providing
the terminal bound. Section~\ref{sec:plain} proves a mean-square estimate
for products of two plain polynomials, again with permitted prime factors
(Lemma~\ref{lem:plain}). Ordinary Hecke reflection reduces the length
ranges at each stage; the recursive step uses two finite Poisson
transformations.

The plain estimate retains different row families according to whether
prime factors are present. With them, it excludes rows whose primitive
inducing character belongs to the fixed finite group generated by the
target and the auxiliary ray characters; without them, it excludes only
principal inducing characters. Transformed rows in that finite family
are treated separately within the induction. Direct volume bounds handle
uncentered products. For a difference of two products with the same two
profiles and equal products of scales, the common main terms cancel.
This cancellation is used only for that centered expression, not for every
row in the finite inducing-character family.

Section~\ref{sec:conclusion} combines the resulting row counts with the
compensated expansion, its local errors, the contour tails, and the
principal residue. After normalization, their bound verifies the second
comparison in the continuation criterion at $7/8$.

\paragraph{Transfer to Dirichlet $L$-functions.}
The final transfer is the same at both boundaries. Composing a Dirichlet
character $\chi$ with the ideal norm gives a Hecke character over $F$.
Away from finitely many Euler factors, quadratic factorization writes its
$L$-function as the product of the Dirichlet $L$-functions attached to $\chi$
and to $\chi\chi_{-3}$, where $\chi_{-3}$ is the quadratic character of
$F/\mathbb Q$. The omitted factors are nonzero in $\Re s>0$, and the
principal pole is handled separately. Thus each Hecke half-plane transfers
to all Dirichlet characters, completing the two stages.

\subsection*{Least nonresidues and square roots over prime fields}

The fixed Dirichlet zero-free half-plane also makes the following
classical consequences unconditional.

\begin{corollary}\label{cor:nonresidues-square-roots}
There are absolute constants $A,C>0$ such that, for every odd prime $p$,
the least positive quadratic nonresidue $n(p)$ satisfies
\[
 n(p)\le C(\log p)^A.
\]
In particular, $n(p)\ll_\delta p^\delta$ for every $\delta>0$, proving
Vinogradov's least quadratic nonresidue conjecture
\cite[Conjecture 1.1]{TaoNonresidues}.
Given an odd prime $p$ and $a\in\mathbb F_p$ in binary representation,
there is a deterministic algorithm, with running time polynomial in
$\log p$, that returns a square root of $a$ or reports that none exists.
\end{corollary}

\begin{proof}
By Theorem~\ref{thm:main} and the functional equation, the nontrivial
zeros of every primitive Dirichlet $L$-function lie in
$1/8\le\Re s\le7/8$. They therefore lie in the strictly larger strip
$1/16<\Re s<15/16$. This supplies the weak-GRH hypothesis of
Bhargava, Ivanyos, Mittal, and Saxena
\cite[Conjecture 6.3 and Theorem 6.7]{BIMS} with $\epsilon=7/16$.
Their bound for the least quadratic nonresidue gives the asserted
inequality, for example with $A=32$; no optimization is needed here.
The assertion for each $\delta>0$ follows because every fixed power
of $\log p$ is $O_\delta(p^\delta)$.

For the algorithm, first handle $a=0$ and test a nonzero $a$ by Euler's
criterion. If $a$ is a square, scan $2,3,\ldots$ until a quadratic
nonresidue is found, testing each candidate by its Legendre symbol.
The bound just proved makes this a polynomial-time deterministic
search; its stopping rule does not require knowing $C$.
Use the resulting nonresidue in the Tonelli--Shanks algorithm
\cite[Section 2.9, Algorithm 3 and Lemma 2.9.5]{Galbraith}.
All remaining steps are deterministic and polynomial in $\log p$.
Writing $p-1=2^e u$ with $u$ odd only requires removing factors of two,
not factoring $u$.
\end{proof}

\section{From a common signal to a zero-free half-plane}
\label{sec:analytic-reduction}

Both parts of the proof use the same analytic principle. A character sum
is bounded directly and is also compared with a Mellin integral containing
the reciprocal of a target $L$-function. A power saving in both comparisons
then continues that reciprocal across the rightmost possible zeros. We prove
this principle for a variable boundary, so that it can be applied at
$11/12$ and at $7/8$ without repeating the argument.

Throughout the paper, $F=\mathbb Q(\sqrt{-3})$. It suffices to consider
primitive target characters. Indeed, a character induced from a primitive
character $\eta$ has $L$-function differing from $L_F(s,\eta)$ by finitely
many factors $1-\eta(\mathfrak p)N\mathfrak p^{-s}$, all nonzero for
$\Re s>0$. Define
\begin{equation}
 \beta_* = \sup\left(\{1/2\}\cup
 \left\{\Re\rho:\begin{array}{l}
  1/2\le\Re\rho\le1,\quad L_F(\rho,\eta)=0\\
  \text{for some primitive finite-order Hecke character }\eta
 \end{array}\right\}\right).
 \label{old-eq:1.1b}
\end{equation}
Poles are not included. Absolute convergence of the Euler product excludes
zeros in $\Re s>1$, so $1/2\le\beta_*\le1$.

For a finite set $\mathcal S$ of prime ideals, a superscript $\mathcal S$
means that the corresponding Euler factors have been deleted:
\[
 L_F^{\mathcal S}(s,\eta)
 =L_F(s,\eta)\prod_{\mathfrak p\in\mathcal S}
       (1-\eta(\mathfrak p)N\mathfrak p^{-s}).
\]
The local value is zero at a ramified prime. Each displayed factor is nonzero
for $\Re s>0$, so deleting finitely many factors neither creates nor removes
a zero there.

\begin{proposition}[Continuation from a common signal]
\label{lem:continuation-criterion}
Fix $\sigma_0\in(1/2,1)$ and suppose $\Delta_0=\beta_*-\sigma_0>0$.
Let $C(s)=s+c$, where $c\in\mathbb R$ is fixed. Suppose that numbers
$\omega,\sigma$ satisfying
\[
 0<\omega<\Delta_0,\qquad \sigma>0
\]
can be chosen independently of the target character. For every primitive
finite-order Hecke character $\eta$, suppose there exist a finite set
$\mathcal S$, a holomorphic function $H_\eta$ on $\Re s>\sigma_0$, and a
function $J_\eta(Z)$ defined for all sufficiently large real $Z$, such that
\begin{equation}
 \sup_{\Re s>\sigma_0}|H_\eta(s)-1|\le\frac12.
 \label{eq:signal-euler-contract}
\end{equation}
For $Z>0$, define
\begin{equation}
 f_\eta(Z)=\frac1{2\pi i}\int_{\Re s=2}
 Z^{C(s)}e^{(s-5/6)^2}
 \frac{H_\eta(s)}{L_F^{\mathcal S}(s,\eta)}\,ds.
 \label{eq:common-residue}
\end{equation}
Assume that, as $Z\to\infty$,
\begin{align}
 |J_\eta(Z)|&\ll_\eta Z^{C(\sigma_0)+\omega},
       \label{eq:low-probe-contract}\\
 |J_\eta(Z)-f_\eta(Z)|&\ll_\eta Z^{C(\beta_*)-\sigma}.
       \label{eq:high-probe-contract}
\end{align}
The implied constants, lower thresholds, excluded set, and function $H_\eta$
may depend on $\eta$. Then these assumptions contradict $\beta_*>\sigma_0$.
\end{proposition}

\begin{proof}
Put
\[
 \epsilon_*:=\min\{\Delta_0-\omega,\sigma\}>0.
\]
Because $C$ has slope one, the triangle inequality gives
\begin{equation}
 |f_\eta(Z)|\ll_\eta Z^{C(\beta_*)-\epsilon_*}
 \qquad (Z\ge Z_{0,\eta}).
 \label{eq:uniform-saving}
\end{equation}
Moreover $\epsilon_*\le\Delta_0-\omega<\Delta_0$, so
$\beta_*-\epsilon_*>\sigma_0$.

We also need control as $Z\downarrow0$. By
Equation~\eqref{eq:signal-euler-contract}, $H_\eta$ is bounded on
$\Re s>\sigma_0$. The reciprocal Euler product is absolutely and uniformly
bounded on $\Re s\ge2$. On every fixed strip $2\le\Re s\le B$, the Gaussian
in Equation~\eqref{eq:common-residue} is $O_B(e^{-(\Im s)^2})$. Cauchy's
theorem on rectangles therefore moves the contour to any fixed $B>2$, with
horizontal integrals tending to zero. Hence
\[
 |f_\eta(Z)|\ll_{\eta,B}Z^{B+c}\qquad(0<Z\le1).
\]
Since $B$ is arbitrary, the signal has arbitrarily rapid power decay at zero.

This bound and Equation~\eqref{eq:uniform-saving} show that
\[
 F_\eta(s):=\int_0^\infty f_\eta(Z)Z^{-C(s)}\,\frac{dZ}{Z}
\]
converges locally uniformly on $\Re s>\beta_*-\epsilon_*$. On each compact
subset, choose $B$ larger than all occurring real parts for the integral near
zero, and use Equation~\eqref{eq:uniform-saving} near infinity. Thus $F_\eta$
is holomorphic on that half-plane.

We identify this Mellin transform without moving a contour across a zero.
Set
\[
 A_\eta(s)=e^{(s-5/6)^2}
             \frac{H_\eta(s)}{L_F^{\mathcal S}(s,\eta)}
 \qquad(\Re s>1).
\]
The function $A_\eta(2+it)$ is continuous and integrable in $t$. Writing
$Z=e^u$, Equation~\eqref{eq:common-residue} becomes
\[
 e^{-(2+c)u}f_\eta(e^u)
   =\frac1{2\pi}\int_{\mathbb R}e^{itu}A_\eta(2+it)\,dt.
\]
The left side is integrable in $u$ by the two endpoint estimates. Ordinary
Fourier inversion therefore gives $F_\eta(2+it)=A_\eta(2+it)$ for every real
$t$. Both sides are holomorphic on $\Re s>1$, so the identity theorem gives
\[
 F_\eta(s)=e^{(s-5/6)^2}
             \frac{H_\eta(s)}{L_F^{\mathcal S}(s,\eta)}
 \qquad(\Re s>1).
\]

Equation~\eqref{eq:signal-euler-contract} implies $|H_\eta(s)|\ge1/2$ on
$\Re s>\sigma_0$. Consequently
\[
 e^{-(s-5/6)^2}\frac{F_\eta(s)}{H_\eta(s)}
 \qquad(\Re s>\beta_*-\epsilon_*)
\]
is a holomorphic continuation of $1/L_F^{\mathcal S}(s,\eta)$.
The number $\epsilon_*$ was chosen independently of $\eta$. By the definition
of $\beta_*$, some target has a zero $\rho$ with
$\Re\rho>\beta_*-\epsilon_*$. Its reciprocal has a pole at $\rho$, and the
deleted factors are nonzero there. This contradicts the continuation.
\end{proof}

The high estimate in Proposition~\ref{lem:continuation-criterion} is measured
relative to $C(\beta_*)$, not to $C(\sigma_0)$. This distinction matters when
a row contribution reaches the low scale but still has a power saving relative
to the hypothetical rightmost zero. Part~\ref{part:eleven-twelfths} uses
\[
 \sigma_0=\frac{11}{12},\qquad C(s)=s-\frac23,
 \qquad C(\sigma_0)=\frac14,
\]
whereas Part~\ref{part:seven-eighths} uses
\[
 \sigma_0=\frac78,\qquad C(s)=s-\frac{11}{16},
 \qquad C(\sigma_0)=\frac3{16}.
\]
Only the positive power margins must be uniform in the target. Fixed-character
constants, excluded sets, and sufficiently large lower thresholds may depend
on it in both applications.

\clearpage
\part{The quasi-Riemann hypothesis}
\label{part:eleven-twelfths}

\section{A first zero-free half-plane}
\label{sec:first-stage}

We first prove a fixed zero-free half-plane using the basic completed
cubic-theta sum. This isolates the mechanism that excludes zeros before the
additional estimates needed for the sharper boundary are introduced.

\begin{theorem}[The $11/12$ half-plane]
\label{thm:eleven-twelfths}
Every finite-order Hecke $L$-function over $F=\mathbb Q(\sqrt{-3})$ has no
zero in $\Re s>11/12$. The same holds for every Dirichlet $L$-function,
including $\zeta(s)$. A pole at $s=1$ for a principal character is allowed.
\end{theorem}

For the Hecke assertion, suppose for contradiction that
\begin{equation}
 \Delta_1:=\beta_*-\frac{11}{12}>0.
 \label{eq:first-stage-gap}
\end{equation}
The proof will construct one character sum with two exact representations.
Cubic-theta reflection bounds the sum after an elementary additive
large-sieve estimate. Poisson summation expresses the same sum as a principal
Mellin signal and a family of remaining character rows. A zero detector and
the sextic large sieve control those rows. These estimates verify
Proposition~\ref{lem:continuation-criterion} with boundary $11/12$.

The fixed family in Equation~\eqref{old-eq:1.1b} is essential even when the
desired consequence concerns $\zeta$: the Poisson representation introduces
finite-order Hecke twists of the target. Proving the family-wide assertion
also supplies the precise input used at the beginning of
Part~\ref{part:seven-eighths}.

\section{Arithmetic and analytic preliminaries}\label{sec:arithmetic}
